\documentclass[10pt,a4paper]{amsart}
\usepackage[margin=19mm]{geometry}
\usepackage{amsmath,amssymb,mathtools}
\usepackage[scr=rsfs,cal=euler]{mathalfa}
\usepackage{xfrac}
\usepackage[unicode,hidelinks,bookmarks=false]{hyperref}
\newcommand{\ie}{i.e.\ }
\newcommand{\cf}{cf.\ }
\newcommand{\resp}{respectively\ }
\newcommand{\Subsection}{\subsection}
\providecommand{\nobreakdash}{\nobreak\hbox{-}}
\setlength{\emergencystretch}{2em}
\multlinegap0pt
\usepackage{bm}
\usepackage{amssymb}
\usepackage{mathtools}
\usepackage{xcolor}
\usepackage{enumerate}
\newtheorem{lemma}{Lemma}
\newcommand{\R}{\mathbb{R}}
\newcommand{\bH}{\mathbb{H}}
\newcommand{\bHS}{\mathrm{HS}}
\newcommand{\norm}[1]{\left\lVert#1\right\rVert}
\renewcommand{\phi}{\varphi}
\DeclareMathOperator{\vol}{Vol}
\title{Quantum Mixing for Schr\"odinger eigenfunctions in Benjamini-Schramm limit}
\author{Kai Hippi, F\'{e}lix Lequen, S{\o}ren Mikkelsen, Tuomas Sahlsten, Henrik Uebersch\"{a}r}
\date{Source: arXiv:2604.21582v2 (CC BY 4.0)}
\begin{document}
\maketitle

\noindent\emph{Source and licence.} Quantum Mixing for Schr\"odinger eigenfunctions in Benjamini-Schramm limit. Version arXiv:2604.21582v2 (CC BY 4.0), distributed by arXiv under the Creative Commons Attribution 4.0 International licence, http://creativecommons.org/licenses/by/4.0/, as recorded in the arXiv licence metadata for this exact version and shown on the abstract page https://arxiv.org/abs/2604.21582v2. Full licence notice: \texttt{licenses/arxiv-2604.21582-CC-BY-4.0.txt}.\par\smallskip

\noindent\textit{Editorial scope.} Counted unit: the article's lemma LE:free\_HS\_est\_exponential\_mixing (source0.tex lines 1582--1604). The article's complete original proof of it is delivered with it (source0.tex lines 1606--1719, one \texttt{proof} environment of 114 lines). No mathematics is written, repaired or re-derived: the statement and the proof are the article's own lines, in the article's own notation, and no retained line is substituted or re-flowed.  Retained uncounted as the source-local context the proof needs: lines 1038--1047; lines 1128--1134; lines 1346--1356; lines 1524--1526; lines 1528--1539.  Nothing was omitted from any retained span, so the package carries no omission record.  No retained line contains a citation, so the wrapper carries no bibliography and no bibliography entry.  No retained line contains a figure environment, an image-inclusion command or drawing code, so no image is delivered and the package declares no asset dependency.  Editorial formatting only, in addition: an AMS \texttt{amsart} preamble that restates the article's own declarations and macros used by the retained lines, namely the environments \texttt{lemma} on separate counters, together with the standard packages those lines need.  The retained environments print in this document as Lemma 1 in order of appearance, whereas the article prints the same items with the numbering of its own full document.  The complete native source of the article as distributed in the arXiv e-print for this exact version is bundled unchanged as \texttt{sources/199033/source0.tex}.
\begin{lemma}\label{LE:kernel_of_unper_pro}
    Let $X=\Gamma\setminus \bH$, where $\Gamma$ be a strictly hyperbolic Fuchsian group, and let $-\Delta_X$ be the Laplace-Beltrami operator. Set
    \begin{equation*}
        P_0(t) = h(t,-\Delta_X).
    \end{equation*}
    Then the integral kernel of $P_0(t)$ is given by
    \begin{equation*}
        K^\Gamma_{t}(x,y)= \sum_{\gamma \in \Gamma} A(t,d(x,\gamma y)) =\frac{1}{2\sqrt{2} \pi} \sum_{\gamma \in \Gamma} \frac{\boldsymbol{1}(t> d(x,\gamma y))}{\sqrt{\cosh(t) - \cosh(d(x,\gamma y))}}.
    \end{equation*}
\end{lemma}

\begin{lemma}\label{LE:Two_abel_kernel_int_est}
Let $z,z'\in\bH$ and $t,t'\in \R_{+}$. Then we have the following
\begin{equation*}
    \int_{\bH} A(t,d(z,x))A(t',d(z',x)) \, d x \leq \frac{ \boldsymbol{1}_{[0,t+t']}(d(z,z'))}{\sqrt{\sinh (\max(|t-t'|,d(z,z')))}},
\end{equation*}
where $A(\tau,r)$ is the Abel kernel introduced in Lemma~\ref{LE:kernel_of_unper_pro}.
\end{lemma}

\begin{lemma}\label{LE:Potential_restriction_1}
Let $X=\Gamma\setminus \bH$, where $\Gamma$ is a strictly hyperbolic Fuchsian group, and let $\tau\in\R$. Suppose $V\in L^\infty(X)$ and let $P_0(t)$ and $\widetilde{P}_0(t,\tau)$ be the free and modified free wave propagator associated with $-\Delta_X$. Then for all $T\geq 1$ and $a\in L^\infty(X)$ we have for any $t,t'\in \R_{+}$ such that $\max(t,t')\leq T$ the estimates
%
\begin{equation*}
       \big\lVert  P_0(t) a \boldsymbol{1}_{X( > 2T)} P_0(t') V  \big\rVert_{\bHS}^2 \leq  4\pi (t+t') \norm{a}_{L^\infty(X)}^2 \norm{V}_{L^2(X)}^2
\end{equation*}
and
\begin{equation*}
       \big\lVert  \widetilde{P}_0(t,\tau) a \boldsymbol{1}_{X( > 2T)} P_0(t') V  \big\rVert_{\bHS}^2 \leq  4\pi (t+t') \norm{a}_{L^\infty(X)}^2 \norm{V}_{L^2(X)}^2
\end{equation*}
\end{lemma}

\begin{equation}\label{def_of_F}
   F(t,t',d(w,w')) = \int_{\bH} A(t,d(x,w))   A(t',d(x,w'))  \,d x \int_{\bH}    A(t,d(y,w))   A(t',d(y,w'))  \,d y,
\end{equation} 

\begin{lemma}\label{LE:int_est_after_exp_mix}
    Let $\beta$ be a positive constant smaller than 1 and let $0<t'<t$. Then
    \begin{equation}\label{LEEQ:int_est_after_exp_mix_1}
       F(t,t',r)\leq \frac{\boldsymbol{1}_{[0,t+t'](r)}}{\sinh(\max(t-t',r))},
    \end{equation}
    and moreover this implies that for $\kappa>1$
    \begin{equation}\label{LEEQ:int_est_after_exp_mix_2}
        \int_{0}^{t+t'} \sinh(r) (1+r)^{-\kappa}   F(t,t',r)\, d r \leq \begin{cases} \frac{\kappa}{\kappa-1}(1+t-t')^{1-\kappa}, & 1 < \kappa < 2 \\ 
        2(1+t-t')^{-1},& \kappa \geq 2.
        \end{cases}
    \end{equation}
\end{lemma}

\begin{lemma}\label{LE:free_HS_est_exponential_mixing}
Let $X=\Gamma\setminus \bH$, where $\Gamma$ be a strictly hyperbolic Fuchsian group, and let $\tau\in\R$. Let $P_0(t)$ and $\widetilde{P}_0(t,\tau)$ be the free and modified free wave propagator associated with $-\Delta_X$. Moreover, suppose that $a\in L^\infty(X)$ with zero average and that there exist $C_1>0$ and $\kappa>1$ such that  
\begin{equation}\label{LEEQ:free_HS_est_exponential_mixing_1}
  \big|\langle a \circ \phi_r, a \rangle \big| \leq C_1 (t+1)^{-\kappa} \lVert a \rVert_{L^2(X)}^2,\quad r \geq 0.  
\end{equation}
%
Then for all $T\geq 1$, we have the estimates
%
\begin{equation*}
       \Big\lVert \int_0^T  P_0(t) a \boldsymbol{1}_{X( > 2T)} P_0(t) \,d t  \Big\rVert_{\bHS}^2 \leq g_\kappa(T) \norm{a}^2_{L^2(X)}+  4\pi T^3 \norm{a}_{L^\infty(X)}^2 \vol(X(\leq 2T)),
\end{equation*}
and 
\begin{equation*}
       \Big\lVert \int_0^T  \widetilde{P}_0(t,\tau) a \boldsymbol{1}_{X( > 2T)} P_0(t) \,d t  \Big\rVert_{\bHS}^2 \leq g_\kappa(T) \norm{a}^2_{L^2(X)}+  4\pi T^3 \norm{a}_{L^\infty(X)}^2 \vol(X(\leq 2T)),
\end{equation*}
where 
\begin{equation*}
    g_\kappa(T) = \begin{cases}
        \frac{\kappa C_1}{\kappa -1} T^{3-\kappa} & \text{if $1<\kappa<2$,}\\
        2 C_1(T+1)\ln (T+1) & \text{if $\kappa\geq 2$.}
    \end{cases}        
\end{equation*}
\end{lemma}

\begin{proof}
We will here give the proof for the second estimate. The first is obtained by an analogous argument that omits the factor $\cos(\tau t)$. Letting $K_{t}^\Gamma(x,y)$ be the integral kernel of $P_0(t)$, we have
\begin{equation*}
   \Big\lVert \int_0^T  \widetilde{P}_0(t,\tau) a \boldsymbol{1}_{X( > 2T)} P_0(t) \,d t  \Big\rVert_{\bHS}^2 = \int_D \int_D \big|\int_0^T \cos(\tau t) \int_{D( > 2T)} K_t^\Gamma(x,w) a(w) K_{t}^\Gamma(w,y) \, d w d t  \big|^2 \, d yd x.
\end{equation*}
Recall that
\begin{equation*}
        K^\Gamma_{t}(x,y)= \sum_{\gamma \in \Gamma} A(t,d(x,\gamma y)).
\end{equation*}
By arguing as in the proof of Lemma~\ref{LE:Potential_restriction_1} we obtain after expanding the square the following identity
\begin{equation*}
    \begin{aligned}
         \MoveEqLeft \Big\lVert \int_0^T  P_0(t) a \boldsymbol{1}_{X( > 2T)} P_0(t) \,d t  \Big\rVert_{\bHS}^2 
         \\
         ={}&  \Big| \int_D \int_D \big| \int_0^T \cos(\tau t) \int_{D( > 2T)} \sum_{\gamma_1,\gamma_2\in\Gamma} A(t,d(\gamma_1 x,w)) a(w) A(t,d((w, \gamma_2 y)) \, d w  d t \big|^2 \, d yd x \Big|
         \\
         ={}&\Big| \int_0^T\int_0^T \cos(\tau t) \cos(\tau t') \int_{D( > 2T)}  \int_{B(w,2T)\cap \bH( > 2T) }  a(w) \overline{a(w')}  F(t,t',d(w,w'))d w' d w      d t'd t\Big|,
    \end{aligned}
\end{equation*}
where we view $a$ as a periodic function on $\bH$,  $ \bH(>2T) = \bigcup_{\gamma \in \Gamma} \gamma D(>2T) $, and the function $F(t,t',(d(w,w'))$ is as defined in \eqref{def_of_F}. Moreover, the restriction of the integration domain of $w'$ follows from the support properties of the Abel kernels. We then notice that due to the support properties of $F$ we have that 
\begin{equation*}
    \begin{aligned}
         \MoveEqLeft  \int_{D( > 2T)}  \int_{B(w,2T)\cap \bH( > 2T) }  a(w) \overline{a(w')}  F(t,t',d(w,w')) dx dyd w' d w  
         \\
         &=  \int_{D( > 2T)}  \int_{B(w,2T)\cap D( > 2T) }  a(w) \overline{a(w')} \sum_{\gamma \in\Gamma}  F(t,t',d(w, \gamma w')) dx dyd w' d w 
         \\
         &=\int_{D( > 2T)}  \int_{B(w,2T)\cap D( > 2T) }  a(w) \overline{a(w')}  F(t,t',d(w, w')) dx dyd w' d w, 
    \end{aligned}
\end{equation*}
where an argument analogous to the one used in the proof of Lemma~\ref{LE:Potential_restriction_1} gives us that only the identity element contributes. Using this we obtain that 
\begin{equation}\label{EQ:free_HS_est_exponential_mixing_1}
    \begin{aligned}
         \MoveEqLeft \Big\lVert \int_0^T  P_0(t) a \boldsymbol{1}_{X( > 2T)} P_0(t) \,d t  \Big\rVert_{\bHS}^2 
         \\
         ={}&\Big| \int_0^T\int_0^T \cos(\tau t) \cos(\tau t') \int_{D( > 2T)}  \int_{B(w,2T)\cap D( > 2T) }  a(w) \overline{a(w')}  F(t,t',d(w,w'))d w' d w      d t'd t\Big|,
    \end{aligned}
\end{equation}
By adding and subtracting the two terms, 
\begin{equation*}
     \int_0^T\int_0^T \cos(\tau t) \cos(\tau t') \int_{D( > 2T)}  \int_{B(w,2T)\cap D( \leq 2T) } a(w) \overline{a(w')} F(t,t',d(w,w'))  \,d w' d w      d t'd t
\end{equation*}
and
\begin{equation*}
     \int_0^T\int_0^T \cos(\tau t) \cos(\tau t') \int_{D( \leq 2T)}  \int_{B(w,2T)} a(w) \overline{a(w')}  F(t,t',d(w,w')) \,d w' d w      d t'd t.
\end{equation*}
We have that
\begin{equation}\label{EQ:free_HS_est_exponential_mixing_3}
    \begin{aligned}
         \MoveEqLeft \Big| \int_0^T\int_0^T \cos(\tau t) \cos(\tau t') \int_{D( > 2T)}  \int_{B(w,2T)\cap \bH( > 2T) } a(w) \overline{a(w')} F(t,t',d(w,w'))  \,d w' d w      d t'd t\Big|
         \\
         \leq {}&  \int_0^T\int_0^T \Big| \int_{D}  \int_{B(w,2T) } a(w) \overline{a(w')}  F(t,t',d(w,w')) \,d w' d w   \Big|   d t'd t
         \\
         &+ \int_0^T\int_0^T \Big| \int_{D( > 2T)}  \int_{B(w,2T)\cap D( \leq 2T) } a(w) \overline{a(w')} F(t,t',d(w,w'))  \,d w' d w    \Big|  d t'd t 
         \\
         &+ \int_0^T\int_0^T \Big| \int_{D( \leq 2T)}  \int_{B(w,2T)} a(w) \overline{a(w')}  F(t,t',d(w,w'))  \,d w' d w  \Big|    d t'd t 
         \\
         ={}& \mathcal{I}_1(a,T) +\mathcal{I}_2(a,T)+\mathcal{I}_3(a,T),
    \end{aligned}
\end{equation}
where we have used $|\cos(\tau t)|\leq 1$. We will first estimate $\mathcal{I}_2(a,T)$ and $\mathcal{I}_3(a,T)$. To do so, we notice that by applying Lemma~\ref{LE:Two_abel_kernel_int_est} to the integrals of the variables $x$ and $y$ in the definition of the function $ F(t,t',d(w,w'))$, we obtain the estimate 
\begin{equation}\label{EQ:free_HS_est_exponential_mixing_4}
    F(t,t',d(w,w')) \leq  \frac{\boldsymbol{1}_{[0,t'+t]}(d(w,w'))}{\sinh(\max(|t-t'|,d(w,w'))}.
\end{equation}
Moreover, since $ F(t,t',d(w,w'))$ is symmetric in $w$ and $w'$, we obtain the following estimate
\begin{equation}\label{EQ:free_HS_est_exponential_mixing_5}
    \begin{aligned}
        \mathcal{I}_2(a,T)+\mathcal{I}_3(a,T)
        &\leq 
         2\norm{a}^2_{L^\infty(X)}\int_0^T\int_0^T \int_{D} \int_{D( \leq 2T)} \frac{\boldsymbol{1}_{[0,t'+t]}(d(w,w'))}{\sinh(\max(|t-t'|,d(w,w'))}   \, d w  d w'   d t'd t 
         \\
         &\leq 4\pi T^3\norm{a}^2_{L^\infty(X)} \vol(X(\leq 2T)),
    \end{aligned}
\end{equation}
where we first evaluated the integral in $w'$ by switching to polar coordinates centered at $w$ and then the remaining integrals were evaluated. 
To estimate $\mathcal{I}_1(a,T)$ we first notice that $F(t,t',d(w,w'))$ is symmetric in $t$ and $t'$. Therefore, we obtain the following
\begin{equation*}
    \begin{aligned}
        \mathcal{I}_1(a,T)
        \leq 2  \int_0^T\int_0^t \Big| \int_{D}  \int_{B(w,2T) } a(w) \overline{a(w')}  F(t,t',d(w,w'))  \,d w' d w   \Big|   d t'd t
    \end{aligned}
\end{equation*}
Next we write $w'$ in polar coordinates around $w$ such that $w' = \pi_1(\varphi_r(w,\theta))$, this gives us 
\begin{equation*}
    \begin{aligned}
        \mathcal{I}_1(a,T)
        &\leq 2  \int_0^T\int_0^t \Big| \int_{D} \int_{0}^{2T}  \int_{\mathbb{S}^1 } a(w) \overline{a(\pi_1(\varphi_r(w,\theta))}  \sinh(r) F(t,t',r)  \,d \theta d r d w   \Big|   d t'd t
        \\
        &\leq 2 \int_0^T\int_0^t \int_{0}^{2T} \sinh(r)F(t,t',r)  \Big| \int_{D}   \int_{\mathbb{S}^1 } a(w) \overline{a(\pi_1(\varphi_r(w,\theta))}  \,d \theta d w   \Big| d r   d t'd t
        \\
        &= 2 \int_0^T\int_0^t \int_{0}^{2T} \sinh(r)F(t,t',r)  \big| \langle a, a\circ \varphi_r\rangle_{L^2(D\times \mathbb{S}^1)}   \big| d r   d t'd t.
    \end{aligned}
\end{equation*}
From our assumptions \eqref{LEEQ:free_HS_est_exponential_mixing_1} we have that there exist a $C_1>0$ and $\kappa>1$ such that 
\begin{equation*}
  \big|\langle a \circ \phi_r, a \rangle \big| \leq C_1 (t+1)^{-\kappa} \lVert a \rVert_{L^2(X)}^2,\quad r \geq 0.  
\end{equation*}
Hence, using Lemma~\ref{LE:int_est_after_exp_mix} we obtain the following estimate
\begin{equation}\label{EQ:free_HS_est_exponential_mixing_7}
    \begin{aligned}
        \mathcal{I}_1(a,T)
        &\leq C_1 \max\big(\tfrac{\kappa}{\kappa-1},2\big)\norm{a}^2_{L^2(X)} \int_0^T\int_0^t (1+t-t')^{\max (1-\kappa,-1)}   d t'd t
        \\
        & \leq  g_\kappa(T) \norm{a}^2_{L^2(X)},
    \end{aligned}
\end{equation}
where
\begin{equation*}
    g_\kappa(T) = \begin{cases}
        \frac{\kappa 2^{4-\kappa} C_1}{(\kappa -1)(2-\kappa)} T^{3-\kappa} & \text{if $1<\kappa<2$,}\\
        4 C_1(T+1)\ln (T+1) & \text{if $\kappa\geq 2$.}
    \end{cases}        
\end{equation*}
To finalise the proof, we combine the estimates and identities in \eqref{EQ:free_HS_est_exponential_mixing_1}, \eqref{EQ:free_HS_est_exponential_mixing_3}, \eqref{EQ:free_HS_est_exponential_mixing_5}, and \eqref{EQ:free_HS_est_exponential_mixing_7} to obtain the stated estimate. This concludes the proof.    
\end{proof}

\end{document}
